% =====================================================================
% QUESTIONÁRIO DE LEITURA BETA — Molambudos: O Diário do Paciente 1.260
% Autor: Marcelo Claro
% Uso: entregar ao leitor beta após a leitura integral do manuscrito.
% Compilar: pdflatex questionario_leitura_beta.tex
% =====================================================================
\documentclass[11pt,a4paper]{article}
\usepackage[top=2.2cm,bottom=2.2cm,left=2.4cm,right=2.4cm]{geometry}
\usepackage[T1]{fontenc}
\usepackage[utf8]{inputenc}
\usepackage[brazilian]{babel}
\usepackage{enumitem}
\usepackage{booktabs}
\usepackage{array}
\usepackage{xcolor}
\usepackage{hyperref}
\hypersetup{colorlinks=true,linkcolor=blue!60!black,urlcolor=blue!60!black}

\setlist[enumerate]{leftmargin=*,itemsep=2pt,topsep=4pt}
\setlength{\parindent}{0pt}
\setlength{\parskip}{6pt}

\newcommand{\linha}{\vspace{8pt}\rule{\textwidth}{0.4pt}\vspace{4pt}}
\newcommand{\resp}[1]{\vspace{2pt}\textbf{#1}\vspace{4pt}}

\begin{document}

\begin{center}
  {\LARGE\bfseries QUESTIONÁRIO DE LEITURA BETA}\\[6pt]
  {\large \emph{Molambudos: O Diário do Paciente 1.260}}\\[4pt]
  Marcelo Claro\par
  \vspace{6pt}
  \fbox{\parbox{0.92\textwidth}{\centering
  Este questionário deve ser preenchido \textbf{após} a leitura integral do manuscrito.\\
  Não há respostas certas ou erradas: buscamos sua \textbf{experiência de leitura}, não sua concordância.\\
  Tempo estimado de preenchimento: 30--45 minutos.
  }}
\end{center}

\vspace{10pt}
\section*{Identificação do leitor (opcional)}
\begin{tabular}{@{}p{0.3\textwidth}p{0.62\textwidth}@{}}
Nome ou pseudônimo: & \dotfill\\[4pt]
Data da leitura: & \dotfill\\[4pt]
Como recebeu o manuscrito: & \dotfill\\
\end{tabular}

\section*{A. Perfil do leitor}
\begin{enumerate}[label=A\arabic*.]
\item Faixa etária: \quad ( ) 16--25 \quad ( ) 26--40 \quad ( ) 41--60 \quad ( ) 60+ \quad ( ) prefiro não dizer
\item Quantos livros você lê por ano, em média? \hfill \resp{() 0--2 \quad () 3--6 \quad () 7--12 \quad () 13+}
\item Quais gêneros você costuma ler? (marque todos que se aplicam)
  \begin{itemize}[label=\textbullet]
  \item ( ) Ficção literária \quad ( ) Horror \quad ( ) Realismo mágico \quad ( ) Romance histórico
  \item ( ) Thriller/suspense \quad ( ) Diários/memórias \quad ( ) Não ficção \quad ( ) Outro: \dotfill
  \end{itemize}
\item Você tem formação ou experiência profissional/educacional relacionada a saúde mental?
  \hfill \resp{() Sim \quad () Não \quad () Parcial}
\end{enumerate}

\section*{B. Experiência de leitura}
Responda com escala de \textbf{1 (discordo totalmente)} a \textbf{5 (concordo totalmente)}.

\begin{center}
\renewcommand{\arraystretch}{1.5}
\begin{tabular}{@{}p{0.62\textwidth}>{\centering\arraybackslash}p{0.8cm}>{\centering\arraybackslash}p{0.8cm}>{\centering\arraybackslash}p{0.8cm}>{\centering\arraybackslash}p{0.8cm}>{\centering\arraybackslash}p{0.8cm}@{}}
\toprule
Afirmação & 1 & 2 & 3 & 4 & 5 \\
\midrule
A leitura me prendeu do início ao fim &  &  &  &  &  \\
A estrutura em fragmentos (diários, laudos, anotações) funcionou bem &  &  &  &  &  \\
Em algum momento me senti perdido(a) ou confuso(a) com a ordem dos fragmentos &  &  &  &  &  \\
Eu queria continuar lendo quando precisava parar &  &  &  &  &  \\
O livro me provocou emoções fortes (desconforto, tristeza, suspense, beleza) &  &  &  &  &  \\
O ritmo da Parte 5 (Contaminação) foi adequado &  &  &  &  &  \\
O final (Epílogo e pós-leitura) me pareceu coerente com o restante &  &  &  &  &  \\
\bottomrule
\end{tabular}
\end{center}

\section*{C. Dimensões literárias}
\begin{enumerate}[label=C\arabic*]
\item O que mais funcionou na obra? (pode marcar vários e comentar)
  \begin{itemize}[label=\textbullet]
  \item ( ) Voz dos documentos (laudos, diários, notas)
  \item ( ) Personagens — especialmente o Paciente 1.260 e Lúcia
  \item ( ) Atmosfera/ambiente da Colônia Santa Maria
  \item ( ) Trama e mistério (o que aconteceu, a contaminação)
  \item ( ) Linguagem e estilo
  \item ( ) O projeto gráfico/arquivístico (mapas, rotas de leitura, passaporte)
  \item ( ) Outro: \dotfill
  \end{itemize}
\item Cite \textbf{um momento} que ficou marcado para você (pode ser cena, frase, fragmento):\\[4pt]
  \rule{\textwidth}{0.3pt}\\[2pt]
  \rule{\textwidth}{0.3pt}
\item O que funcionou \textbf{menos}? O que você cortaria ou encurtaria?\\[4pt]
  \rule{\textwidth}{0.3pt}\\[2pt]
  \rule{\textwidth}{0.3pt}
\item Os mapas, grafos e rotas de leitura ajudaram ou atrapalharam a imersão? Comente.\\[4pt]
  \rule{\textwidth}{0.3pt}\\[2pt]
  \rule{\textwidth}{0.3pt}
\end{enumerate}

\section*{D. Ética e responsabilidade (importante para o autor)}
\begin{enumerate}[label=D\arabic*]
\item Você leu o \textbf{AVISO AO LEITOR} (página 5)? \quad \resp{() Sim, antes de começar \quad () Sim, depois \quad () Não li}
\item Se leu: ele foi \textbf{suficiente} e bem posicionado? Comente.\\[4pt]
  \rule{\textwidth}{0.3pt}\\[2pt]
  \rule{\textwidth}{0.3pt}
\item A representação do sofrimento psíquico (diário, laudos, internação) lhe pareceu \textbf{respeitosa}, ou houve algum trecho que lhe pareceu sensacionalista, estigmatizante ou desrespeitoso?\\[4pt]
  \rule{\textwidth}{0.3pt}\\[2pt]
  \rule{\textwidth}{0.3pt}
\item A \textbf{Nota Histórica} (contexto das antigas colônias de isolamento) ajudou a situar a obra? \quad \resp{() Sim \quad () Parcialmente \quad () Não \quad () Não li}
\item Você recomendaria esta obra a alguém que vive ou convive com transtorno mental? Por quê?\\[4pt]
  \rule{\textwidth}{0.3pt}\\[2pt]
  \rule{\textwidth}{0.3pt}
\end{enumerate}

\section*{E. Avaliação global}
\begin{enumerate}[label=E\arabic*]
\item Nota geral (0--10): \hfill \resp{\Large\rule{3cm}{0.3pt}}
\item Os três maiores \textbf{trunfos} da obra:\\[4pt]
  1. \dotfill\\[4pt]
  2. \dotfill\\[4pt]
  3. \dotfill
\item As três maiores \textbf{fragilidades} (se houver):\\[4pt]
  1. \dotfill\\[4pt]
  2. \dotfill\\[4pt]
  3. \dotfill
\item Se pudesse mudar \textbf{uma} coisa na obra, o que seria?\\[4pt]
  \rule{\textwidth}{0.3pt}\\[2pt]
  \rule{\textwidth}{0.3pt}
\item Para quem você indicaria este livro? (público-alvo que você enxerga)\\[4pt]
  \rule{\textwidth}{0.3pt}\\[2pt]
  \rule{\textwidth}{0.3pt}
\end{enumerate}

\section*{F. Encerramento (opcional)}
\begin{enumerate}[label=F\arabic*]
\item Autorizo o autor a usar trechos deste questionário (sem meu nome) em material de divulgação?
  \hfill \resp{() Sim \quad () Não}
\item Aceitaria uma conversa de 15--20 minutos para aprofundar? \quad \resp{() Sim \quad () Não}
  \quad Contato: \dotfill
\item Espaço livre (o que mais quiser dizer):\\[6pt]
  \rule{\textwidth}{0.3pt}\\[2pt]
  \rule{\textwidth}{0.3pt}\\[2pt]
  \rule{\textwidth}{0.3pt}\\[2pt]
  \rule{\textwidth}{0.3pt}
\end{enumerate}

\vspace{12pt}
\begin{center}
\emph{Muito obrigado pela sua leitura atenta. Sua resposta é parte essencial do processo editorial.}\\[4pt]
\textbf{Marcelo Claro} — \href{mailto:contato@marceloclaro.com.br}{contato@marceloclaro.com.br}
\end{center}

\end{document}
